\cotefile{93}
\scanpage{32}
% archive sheet number: 31
\mspage{6}

Let $y$ be in $Y$. Then $\theta(y)$ is in the norm closed convex hull of $\{t(y):t\in T\}$ by [3, Th. 13, p. 422]. Therefore, given any $\varepsilon>0$ there exist $t_1,\ldots,t_n$ in $T$ and $\alpha_i\geqq 0$ $(i=1,\ldots,n)$ such that $\sum\alpha_i=1$ and
\[
\lVert\theta(y)=\sum\alpha_i t_i(y)\rVert<\varepsilon/(M+1),
\]
where $M=\sup\{\lVert t\rVert:t\in T\}$ (which is finite by the uniform boundedness principle.) Then if $t\geqq t_1\unsure{\mathbin{\mathrm{o}}|\circ}\ldots\mathbin{\mathrm{o}}t_n$,
\[
\begin{aligned}
\lVert\theta(y)-t(y)\rVert&=\lVert t[\theta(y)-\sum\alpha_i t_i(y)]\rVert\\
&\leqq M\varepsilon/(M+1)<\varepsilon,
\end{aligned}
\]
so $\displaystyle\lim_{t\in T}\lVert\theta(y)-t(y)\rVert=0$. \hfill Q.E.D.

In order to apply the theorem to the case of $T=\{t_\pi:\pi\in\unsure{\mathcal{P}|P}\}$ acting on $\mathrm{sba}(\unsure{\mathcal{A}|A},X)$, we must show that for every $\mu$ in $\mathrm{sba}(\mathcal{A},X)$ the set
\[
\{t_\pi\mu:\pi\in\mathcal{P}\}
\]
is conditionally weakly compact in $\mathrm{sba}(\mathcal{A},X)$. This follows immediately from Proposition 1 below. The theorem then implies that $\mathrm{sba}(\mathcal{A},X)$ can be decomposed into the direct sum of the countably additive members of $\mathrm{sba}(\mathcal{A},X)$ and the purely finitely additive members according to the following definition.
